\documentclass[../main.tex]{subfiles}

\begin{document}
\subsection{\textit{Learn Stack Layout}}

\begin{code}
\begin{verbatim}
#include <iostream>
\end{verbatim}
\end{code}

{\Script} di atas digunakan untuk membolehkan penggunaan fungsi-fungsi dan objek-objek yang berada di dalam {\header} {\file} \citecode{iostream}. {\Header} {\file} \citecode{iostream} berisi hal-hal yang dibutuhkan untuk membantu proses {\ioi} dan {\ioo}.

\begin{code}
\begin{verbatim}
using namespace std;
\end{verbatim}
\end{code}

{\Script} \citecode{using namespace std;} adalah {\statement} yang digunakan untuk memberikan perintah kepada {\compiler} bahwa file tersebut akan menggunakan seluruh fungsi yang menjadi bagian dari \citecode{namespace std}.

\begin{code}
\begin{verbatim}
const string registers[] = {
  "RDI", "RSI", "RDX", "RCX", "R8 ", "R9 " };
const int register_max = sizeof registers / sizeof *registers;
struct Function {
  string name;
  string arguments[register_max];
  int argument_count;
};
const int stack_max = 9;
struct Stack {
  void *ptr[stack_max];
  int type[stack_max] = { -1 };
  int count = 3;
  bool filled = false;
};
Stack stack;
\end{verbatim}
\end{code}

{\Script} di atas merupakan kumpulan-kumpulan data dan struktur-struktur data yang diperlukan untuk jalan kerjanya program. Struktur data \citecode{Function} digunakan untuk menyimpan masukan fungsi dan struktur data \citecode{Stack} digunakan untuk menyimpan stuck secara keseluruhan.

\begin{code}
\begin{verbatim}
const int STACK_TYPE_VALUE = 0;
const int STACK_TYPE_FUNCTION = 1;
\end{verbatim}
\end{code}

{\Script} di atas digunakan sebagai pengganti struktur data \citecode{enum} dikarenakan tidak termasuk materi praktikum saat ini. Konstanta di atas hanya berguna untuk mempermudah proses pembacaan kode.

\begin{code}
\begin{verbatim}
void topbar() { cout << "\33[H\33[2J";
  cout << " LEARN STACK LAYOUT\n";
  cout << "\n\n\n"; }
\end{verbatim}
\end{code}

{\Script} di atas digunakan untuk menampilkan judul program. Penampilan judul tersebut dilakukan melalui objek \citecode{cout} yang disediakan oleh {\header} {\file} \citecode{iostream} yang merupakan bagian dari \textit{standard library} bahasa pemrograman C++.

\begin{code}
\begin{verbatim}
void pause() {
  cout << "Press any key to continue. . .\n"; cin.get();
}
\end{verbatim}
\end{code}

{\Script} di atas merupakan sebuah fungsi tanpa kembalian yang ditandakan oleh kata kunci \citecode{void}. Fungsi tersebut akan menghentikan pengguna sejenak sampai dideteksi sebuah masukan apapun dari \textit{keyboard}.

\begin{code}
\begin{verbatim}
void input_invalid() {
  cin.clear(); cin.ignore();
  cout << "Masukan tidak valid\n";
}
\end{verbatim}
\end{code}

{\Script} di atas merupakan sebuah fungsi tanpa kembalian yang ditandakan oleh kata kunci \citecode{void}. Fungsi tersebut akan memberitahu pengguna bahwa sebuah masukan tidak valid dan membersihkan masukan tidak valid tersebut dari {\ioi} {\stream}.

\begin{code}
\begin{verbatim}
int nselector(int l, int u, const string& c = ">> ") {
  int n;
  while (true) {
    cout << c;
    cin >> n;
    if (cin.fail() || n < l || n > u) {
      input_invalid();
      continue;
    } break; }
  return n;
}
int nselector(const string& c = "", int i = 0) {
  int n;
  while (true) {
    cout << c << i << " : ";
    cin >> n;
    if (cin.fail()) {
      input_invalid();
      continue;
    } break; }
  return n;
}
\end{verbatim}
\end{code}

{\Script} di atas merupakan sebuah fungsi dengan kembalian berupa \citecode{int} atau bilangan bulat. Fungsi tersebut akan menyaring masukan dari pengguna sesuai kedua parameter fungsi \citecode{int l} dan \citecode{int u}. Masukan yang dimasukan harus berada dalam lingkup kedua parameter tersebut untuk dianggap sebagai masukan yang valid.

\begin{code}
\begin{verbatim}
char cselector(char y='Y', char n='N', string c = ">> ") {
  char ch;
  while (true) {
    cout << c;
    cin >> ch;
    if (ch > '`' && ch < '{') ch -= 32;
    if (cin.fail() || (ch != 'Y' && ch != 'N')) {
      input_invalid();
      continue;
    } break; }
  return ch;
}
\end{verbatim}
\end{code}

{\Script} di atas merupakan sebuah fungsi yang mengembalikan tipe data \citecode{char} sesuai masukan parameter \citecode{char y} dan \citecode{char n}. Bila masukan yang diterima tidak sesuai parameter yang diberikan maka akan diulangi sampai pilihan yang benar dimasukan.

\begin{code}
\begin{verbatim}
void free_stack(Stack *st) {
  for (int i = 0; i < stack.count; ++i) {
      if (st->type[i] == STACK_TYPE_VALUE)
        delete (int *)st->ptr[i];
      else if (st->type[i] == STACK_TYPE_FUNCTION)
        delete (Function *)st->ptr[i]; }
}
\end{verbatim}
\end{code}

{\Script} di atas merupakan sebuah fungsi tanpa nilai kembalian yang dapat dilihat dari tipe data \citecode{void}. Fungsi tersebut menerima masukan parameter {\pointer} ke tipe data \citecode{Stack}. Fungsi tersebut akan membebaskan semua memori yang digunakan oleh argumen masukan.

\begin{code}
\begin{verbatim}
void* push_to_stack(int type) {
  if (type == STACK_TYPE_FUNCTION)
    return (void *)new Function();
  return (void *)new int;
}
\end{verbatim}
\end{code}

{\Script} di atas merupakan sebuah fungsi yang mengembalikan tipe data \citecode{void*}. Tipe data \citecode{void*} digunakan karena tidak diperbolehkan penggunaan \textit{generic}. Fungsi tersebut dapat mengalokasikan memori dengan tipe data \citecode{Function} atau \citecode{int}.

\begin{code}
\begin{verbatim}
void add_data() {
  if (stack.filled) {
    free_stack(&stack);
    stack.filled = false;
    if (stack.count == 6) stack.count = 9;
    if (stack.count == 3) stack.count = 6; }
  cout << "Current Stack Size: " << stack.count << endl;
  for (int i = 0; i < stack.count; ++i) {
    char c0 = cselector('Y', 'N', "Will You Add Function? (Y/N) : ");
    if (c0 == 'Y') {
      stack.type[i] = STACK_TYPE_FUNCTION;
      stack.ptr[i] = push_to_stack(STACK_TYPE_FUNCTION);
      Function *fp = (Function *)stack.ptr[i];
      char c1 = cselector('Y', 'N', "Have Argument(s)? (Y/N): ");
      cout << "Function - " << i + 1 << " To Push : ";
      getline(cin >> ws, fp->name);
      if (c1 == 'Y') {
        fp->argument_count = nselector(1, register_max,
                         "How many arguments ? ");
        for (int j = 0; j < fp->argument_count; ++j) {
          cout << "Argument - " << j + 1 << " : ";
          getline(cin >> ws, fp->arguments[j]);
        }
      } else fp->argument_count = 0;
    } else if (c0 == 'N') {
      stack.type[i] = STACK_TYPE_VALUE;
      stack.ptr[i] = push_to_stack(STACK_TYPE_VALUE);
      *((int *)stack.ptr[i]) = nselector("Nilai - ", i + 1);
    }
  } stack.filled = true;
}
\end{verbatim}
\end{code}

{\Script} di atas merupakan sebuah fungsi tanpa nilai kembalian atau fungsi dengan kembalian bertipe data \citecode{void}. Fungsi tersebut digunakan untuk membantu penambahan data ke \citecode{stack}. Pada saat fungsi dipanggil akan ditanyakan apakah masukan berupa fungsi atau nilai. Bila masukan berupa nilai maka akan diminta nilai bilangan bulat. Apabila masukan berupa fungsi akan ditanyakan apakah fungsi tersebut memiliki argumen dan akan diminta mengisi data-data berkaitan dengan tipe data \citecode{Function}.

\begin{code}
\begin{verbatim}
void debug_mode() {
  while (true) {
    string arg;
    cout << "Debug Command => ";
    getline(cin >> ws, arg);
    if (arg == "exit") return;
    if (arg == "info stack") {
      int i = 0;
      if (stack.type[i] < 0) {
        cout << "Anda seharusnya tidak berada disini...\n";
        continue; }
      while (true) {
        topbar();
        if (i > stack.count - 1) {
          cout << "REGISTER\n";
          cout << "No Registers!\n";
          cout << "STACK\n";
          cout << "No Stacks!\n";
          pause(); return; }
        cout << "REGISTER\n";
        for (int j = 0; j < register_max; ++j) {
          cout << registers[j] << " : ";
          if (stack.type[i] == STACK_TYPE_FUNCTION) {
            Function *fp = (Function *)stack.ptr[i];
            if (j < fp->argument_count) {
              cout << fp->arguments[j] << endl;
              continue; }
          } cout << "0\n"; }
        cout << "STACK\n";
        for (int j = 0; j < stack.count; ++j) {
          if (j == i) cout << "=> ";
          else cout << "   ";
          cout << stack.ptr[j] << "   ";
          if (stack.type[j] == STACK_TYPE_FUNCTION) {
            cout << ((Function *)stack.ptr[j])->name << "\n";
          } else if (stack.type[j] == STACK_TYPE_VALUE) {
            cout << *((int *)stack.ptr[j]) << "\n";
          } }
        cout << "Debug Command => ";
        getline(cin >> ws, arg);
        if (arg == "ni") i++;
        if (arg == "pi") i--;
        if (i < 0) i = stack.count - 1;
      } } } }
\end{verbatim}
\end{code}

{\Script} di atas merupakan fungsi tanpa nilai kembalian yang digunakan untuk menulusuri data-data yang tersimpan dalam \citecode{stack}. Penampilannya akan dipisah menjadi dua bagian yakni \textit{stack} dan register yang akan berubah sesuai nilai yang terkandung pada variabel \citecode{stack}.

\begin{code}
\begin{verbatim}
int main() {
  int state = 0;
  while (true) {
    topbar();
    cout << "1. Add Data\n";
    cout << "2. Debug Mode\n";
    cout << "3. Exit\n";
    state = (nselector(1, 3) - 1);
    switch (state) {
      case 0: {
        add_data();
      } break;
      case 1: {
        debug_mode();
      } break;
      case 2: return 0; } } }
\end{verbatim}
\end{code}

{\Script} di atas merupakan {\entry} {\point} program \textit{Learn Stack Layout}. {\Script} di atas akan dijalankan setiap kali program tersebut dijalankan. Fungsi di atas digunakan untuk mengatur keseluruhan jalan kerja program seperti menu mana yang akan ditampilkan.

\subsection{\textit{Dungeon Crawler RPG}}

\begin{code}
\begin{verbatim}
#include <iostream>
\end{verbatim}
\end{code}

{\Script} di atas digunakan untuk membolehkan penggunaan fungsi-fungsi dan objek-objek yang berada di dalam {\header} {\file} \citecode{iostream}. {\Header} {\file} \citecode{iostream} berisi hal-hal yang dibutuhkan untuk membantu proses {\ioi} dan {\ioo}.

\begin{code}
\begin{verbatim}
using namespace std;
\end{verbatim}
\end{code}

{\Script} \citecode{using namespace std;} adalah {\statement} yang digunakan untuk memberikan perintah kepada {\compiler} bahwa file tersebut akan menggunakan seluruh fungsi yang menjadi bagian dari \citecode{namespace std}.

\begin{code}
\begin{verbatim}
const int GAME_STATE_PLAYING = 0;
const int GAME_STATE_LOSE = 1;
const int GAME_STATE_WIN  = 2;
\end{verbatim}
\end{code}

{\Script} di atas digunakan sebagai pengganti struktur data \citecode{enum} dikarenakan tidak termasuk materi praktikum saat ini. Konstanta di atas hanya berguna untuk mempermudah proses pembacaan kode.

\begin{code}
\begin{verbatim}
struct Monster { string name, sprite;
  int cur_health, max_health, attack, defense, turn; };
struct Item {
  int health, attack, defense;
  string name; };
struct Player { int health, attack, defense;
  int item_id, turn; string name; };
struct Room {
  int monster_id, item_id, prev, next[2]; bool has_branch;
  string name; };
struct Game {
  int state, current_room_id;
  Player players[2];
  Monster monsters[2];
  Item items[2];
  Room rooms[7]; };
\end{verbatim}
\end{code}

{\Script} di atas merupakan kumpulan-kumpulan data dan struktur-struktur data yang diperlukan untuk jalan kerjanya program. Struktur data \citecode{Monster} digunakan untuk menyimpan keadaan \textit{monster}, struktur data \citecode{Item} digunakan untuk menyimpan data \textit{item}, struktur data \citecode{Player} menyimpan data pemain, struktur data \citecode{Room} digunakan untuk menyimpan data ruangan seperti apakah ruangan mengandung \textit{monster} atau \textit{item} dan struktur data \citecode{Game} digunakan untuk menyimpan keseluran perjalanan program.

\begin{code}
\begin{verbatim}
string strpad(int n, char c = ' ') {
  string pad = "";
  for (int i = 0; i < n; ++i) pad += c;
  return pad; }
\end{verbatim}
\end{code}

{\Script} di atas merupakan sebuah fungsi yang digunakan untuk memasukan sebuah karakter sesuai parameter \citecode{char c} sebanyak nilai parameter \citecode{int n}. Fungsi tersebut akan mengembalikan tipe data \citecode{string}.

\begin{code}
\begin{verbatim}
int strlen(const string& text) {
  if (text[0] == '\0') return 0;
  int len = 0;
  while (text[++len] != '\0') {}
  return len; }
\end{verbatim}
\end{code}

{\Script} di atas merupakan sebuah fungsi yang digunakan untuk menemukan panjang sebuah tipe data \citecode{string}. Fungsi tersebut akan mengembalikan nilai bilangan bulat atau \citecode{int} yang berupa panjang \citecode{string} masukan.

\begin{code}
\begin{verbatim}
void strrev(string& text) {
  string s;
  for (int i = strlen(text) - 1; i >= 0; --i)
    s += text[i];
  text = s; }
\end{verbatim}
\end{code}

{\Script} di atas digunakan untuk membalik urutan \citecode{string}. Fungsi tersebut menerima parameter masukan \textit{pass-by} {\reference} yang berarti argumen masukan tersebut akan dimodifikasi oleh fungsi tersebut.

\begin{code}
\begin{verbatim}
string strint(int n) {
  string s = "";
  while (n > 0) {
    s += (char)((n % 10) + 48);
    n /= 10.0f;
  }
  strrev(s); return s; }
\end{verbatim}
\end{code}

{\Script} di atas digunakan untuk mengubah nilai tipe data \citecode{int} menjadi keluaran tipe data \citecode{string}. Fungsi tersebut akan melakukan perulangan pada parameter masukan yang kemudian akan disimpan pada sebuah \citecode{string} yang merupakan kembalian fungsi tersebut.

\begin{code}
\begin{verbatim}
string strup(string text) {
  for (int i = 0; i < strlen(text); ++i)
    if (text[i] > '`' && text[i] < '{')
      text[i] -= 32;
  return text; }
\end{verbatim}
\end{code}

{\Script} di atas digunakan untuk mengubah huruf kecil pada suatu kata menjadi huruf kapital. Setiap karakter pada parameter \citecode{string text} akan dikurangi dengan 32 dikarenakan perbedaan kode ASCII huruf kapital terletak 32 lebih bawah dari huruf kecil.

\begin{code}
\begin{verbatim}
void center_draw(const string& text,
            int width = 51,
            const string& border = "| |") {
  string lpad = strpad(width/2 - strlen(border)
              - (strlen(text)/2));
  string rpad = strpad(45 - (strlen(lpad) + strlen(text)));
  cout << border << lpad << text << rpad << border << endl; }
\end{verbatim}
\end{code}

{\Script} di atas digunakan untuk menggambar teks tepat di tengah sesuai panjang masukan \citecode{int width}. Fungsi tersebut akan menggunakan fungsi \citecode{strpad} untuk memberikan ruang kosong agar teks bisa digambar di tengah.

\begin{code}
\begin{verbatim}
void pause() {
  cout << "Press any key to continue..." << endl;
  string arg; getline(cin, arg); }
\end{verbatim}
\end{code}

{\Script} di atas merupakan sebuah fungsi tanpa kembalian yang ditandakan oleh kata kunci \citecode{void}. Fungsi tersebut akan menghentikan pengguna sejenak sampai dideteksi sebuah masukan apapun dari \textit{keyboard}.

\begin{code}
\begin{verbatim}
void bar_draw() {
  cout << "|+|=============================================|+|\n";}
\end{verbatim}
\end{code}

{\Script} di atas merupakan sebuah fungsi tanpa nilai balik yang digunakan untuk menggambar sebuah garis. Fungsi tersebut memanfaatkan objek \citecode{cout} untuk menggambar baris tersebut.

\begin{code}
\begin{verbatim}
void map_draw(int room_id) {
  cout << "                          (#)\n"
          "                           | \n"
          "                           | \n"
          "                          (#)\n"
          "                           | \n"
          "                           | \n"
          "                          (#)\n"
          "                          / \\_\n"
          "                         /     \\_\n"
          "                        /       (#)\n"
          "                (#)---(#)        |\n"
          "                                /\n"
          "                               |\n"
          "                                \\ \n"
          "                                (@)\n";
\end{verbatim}
\end{code}

{\Script} di atas merupakan potongan dari fungsi \citecode{map\_draw} yang akan digunakan untuk menggambar peta permainan. Penggambaran peta dilakukan dengan cara mengeluarkannya ke objek \citecode{cout}.

\begin{code}
\begin{verbatim}
int line = 4, column = 28;
switch (room_id) {
  case 1: line = 7;
  break; case 2:
  line = 10; break;
  case 3: line = 14;
  column = 24; break;
  case 4: line = 14;
  column = 18; break; case 5: 
  line = 13; column = 34; break; case 6:
  line = 18; column = 34; break; }
cout << "\e7";
cout << "\33[" << line <<";" << column << "H!";
cout << "\e8"; }
\end{verbatim}
\end{code}

{\Script} di atas merupakan potongan bawah dari fungsi \citecode{map\_draw} yang akan digunakan untuk menggambar letak pemain. Letak pemain digambar dengan cara menggerakan kursor konsol ke posisi pemain saat ini menggunakan karakter ANSI.

\begin{code}
\begin{verbatim}
void item_info(const Item& item) {
  if (item.health != 0) 
    center_draw("Health " +
      string{(item.health > 0 ? "Increased By " : "Decreased By ")}
      + strint(item.health));
  if (item.attack != 0)
    center_draw("Attack " +
      string{(item.attack > 0 ? "Increased By " : "Decreased By ")}
      + strint(item.attack));
  if (item.defense != 0)
    center_draw("Defense " +
      string{(item.defense > 0 ? "Increased By " : "Decreased By ")}
      + strint(item.defense)); }
\end{verbatim}
\end{code}

% Placeholder
{\Script} di atas merupakan potongan bawah dari fungsi \citecode{map\_draw} yang akan digunakan untuk menggambar letak pemain. Letak pemain digambar dengan cara menggerakan kursor konsol ke posisi pemain saat ini menggunakan karakter ANSI.
% Placeholder


\end{document}
